% =============================================================
%  数学家 LaTeX 文档共用导言区（中文版）
%  被 pages/<世纪>/<Name>/latex_zh/<Name>_zh.tex 引入
% =============================================================

\documentclass[11pt,a4paper]{article}

% ---- 编码与字体（xelatex）----
\usepackage{fontspec}
\defaultfontfeatures{Ligatures=TeX}

% 西文主字体：Times New Roman / Helvetica / Menlo（macOS）
\IfFontExistsTF{Times New Roman}{\setmainfont{Times New Roman}}{}
\IfFontExistsTF{Helvetica}{\setsansfont{Helvetica}}{}
\IfFontExistsTF{Menlo}{\setmonofont{Menlo}}{}

% ---- 中文支持（xeCJK + 苹方/思源）----
\usepackage{xeCJK}
\xeCJKsetup{CJKecglue={}}  % 中英文之间不强制空格
% 优先使用 macOS 自带的"苹方"，找不到时退回 Heiti / 黑体
\IfFontExistsTF{PingFang SC}{%
  \setCJKmainfont[BoldFont={PingFang SC Semibold},ItalicFont={Kaiti SC}]{PingFang SC}%
}{%
  \IfFontExistsTF{Heiti SC}{\setCJKmainfont{Heiti SC}}{%
    \IfFontExistsTF{STSong}{\setCJKmainfont{STSong}}{%
      \setCJKmainfont{SimSun}%
    }%
  }%
}
\IfFontExistsTF{PingFang SC}{\setCJKsansfont{PingFang SC}}{}
\IfFontExistsTF{Menlo}{\setCJKmonofont{Menlo}}{}
\usepackage{newunicodechar}

% ---- 排版与版面 ----
\usepackage[margin=2.3cm]{geometry}
\usepackage{microtype}
\usepackage{setspace}
\onehalfspacing
\setlength{\parskip}{0.5em}
\setlength{\parindent}{0pt}

% ---- 数学 ----
\usepackage{amsmath,amssymb,amsthm}
\usepackage{mathtools}

% ---- 表格 ----
\usepackage{booktabs}
\usepackage{longtable}
\usepackage{array}
\usepackage{tabularx}
\usepackage{multirow}

% ---- 图像 ----
\usepackage{graphicx}
\usepackage{float}
\graphicspath{{images/}}

% ---- 超链接 ----
\usepackage[hidelinks,breaklinks=true]{hyperref}
\hypersetup{
  colorlinks=true,
  linkcolor=black,
  urlcolor=blue!60!black,
  citecolor=black,
}
\usepackage{url}
\usepackage{xurl}

% ---- 列表 ----
\usepackage{enumitem}
\setlist{topsep=0.2em,itemsep=0.1em,parsep=0.1em}

% ---- pandoc 兼容命令 ----
\providecommand{\tightlist}{%
  \setlength{\itemsep}{0pt}\setlength{\parskip}{0pt}}
\providecommand{\passthrough}[1]{#1}
\providecommand{\pandocbounded}[1]{%
  \makebox[\linewidth][c]{\resizebox{!}{!}{#1}}}

% ---- 段落盒子 ----
\usepackage{framed}
\usepackage{xcolor}

% ---- 页眉页脚 ----
\usepackage{fancyhdr}
\setlength{\headheight}{15pt}
\addtolength{\topmargin}{-3pt}
\pagestyle{fancy}
\fancyhf{}
\fancyhead[L]{\nouppercase{\leftmark}}
\fancyhead[R]{\thepage}
\renewcommand{\headrulewidth}{0.3pt}

\usepackage{titling}

% ---- 语言命令兜底 ----
\AtBeginDocument{%
  \providecommand{\foreignlanguage}[2]{#2}%
  \renewcommand{\foreignlanguage}[2]{#2}%
  \providecommand{\textlang}[2]{#2}%
  \providecommand{\lang}[1]{}%
}

% ---- 元数据命令（由 metadata.tex 填充）----
\providecommand{\MathematicianName}{未知}
\providecommand{\MathematicianBirth}{?}
\providecommand{\MathematicianDeath}{?}
\providecommand{\MathematicianNationality}{}
\providecommand{\MathematicianFields}{}
\providecommand{\MathematicianAwards}{}
\providecommand{\MathematicianDescription}{}
\providecommand{\MathematicianWikipedia}{}

% 信息框（中文标签）
\newcommand{\InfoBox}{%
  \begin{center}
  \begin{tabular}{@{}p{0.28\textwidth}p{0.62\textwidth}@{}}
    \toprule
    \textbf{出生}        & \MathematicianBirth \\
    \textbf{逝世}        & \MathematicianDeath \\
    \textbf{国籍}        & \MathematicianNationality \\
    \textbf{研究领域}    & \MathematicianFields \\
    \textbf{获奖}        & \MathematicianAwards \\
    \textbf{维基百科}    & \url{\MathematicianWikipedia} \\
    \bottomrule
  \end{tabular}
  \end{center}
}
